\documentclass[twocolumn]{ceurart}

%%
%% Rights management.
\copyrightyear{2026}
\copyrightclause{Copyright for this paper by its authors.
  Use permitted under Creative Commons License Attribution 4.0
  International (CC BY 4.0).}

%%
%% Conference information
\conference{CLEF 2026 Working Notes, 22--26 September 2026, Thessaloniki, Greece}

%%
%% Title
\title{Team Fosuai at PAN 2026: Source Retrieval via\\Query-Chunk Provenance Modeling}
\title[mode=sub]{Notebook for the PAN Lab at CLEF 2026}

%%
%% Authors
\author[1]{Yi Wang}[%
  email=wangyi0187@icloud.com,
]
\author[1]{Zhongyuan Han}[%
  email=,
]
\cormark[1]
\author[1]{Chang Liu}[%
  email=,
]
\author[1]{Haojie Cao}[%
  email=,
]

\address[1]{Foshan University, School of Computer Science and Artificial Intelligence,
  Foshan 528000, Guangdong, China}

%%
%% Corresponding author
\cortext[1]{Corresponding author: Zhongyuan Han.}

%%
%% Abstract
\begin{abstract}
The PAN@CLEF 2026 Generative Plagiarism Detection task (Task 4) aims to retrieve
potential source documents for a given suspicious document from a large-scale
corpus. Generative plagiarism poses a novel challenge: LLM-generated suspicious
texts often composite content from multiple distinct sources, causing multi-source
signal dilution when the entire document is treated as a single query. We propose
Query-Chunk Provenance Modeling for source retrieval. The suspicious document is
decomposed into paragraph-level semantic segments; each segment independently
retrieves candidate sources via BM25; and a coverage-weighted voting mechanism
aggregates segment-level results into an explicit provenance graph. Our method
is lightweight---requiring only numpy and the Python standard library---and has
been deployed as a Docker submission on TIRA (team fosuai, software
oscillating-list). On the official PAN26 main test set, our system achieved
RR=0.9350 (ranked 3rd by RR) and improved over the official BM25 baseline on
Recall@10 (0.630 vs.\ 0.577) and nDCG@10 (0.626 vs.\ 0.553). On held-out
development data (7,949 documents, 5,522 queries), paragraph-segmented BM25
reaches R@10=0.9947 (baseline: 0.9804) and MRR=0.9724 (baseline: 0.9174).
\end{abstract}

%%
%% Keywords
\begin{keywords}
  Generative Plagiarism Detection \sep
  Source Retrieval \sep
  Query Segmentation \sep
  Provenance Modeling \sep
  BM25 \sep
  Coverage-Weighted Voting
\end{keywords}

\maketitle


%% ===================================================================
\section{Introduction}

The PAN@CLEF 2026 Generative Plagiarism Detection task (Task~4) addresses a
rapidly emerging challenge: Large Language Models (LLMs) can synthesize,
recombine, and rewrite content from dozens of source documents into a single
suspicious text, where different paragraphs may originate from entirely different
sources~\cite{greinerpetter:2025}. The task consists of two stages:
(1)~Source Retrieval---identify source documents from a large corpus given a
suspicious document; and (2)~Text Alignment---precisely locate plagiarized
passages at character level. This paper focuses on Stage~1.

Existing source retrieval methods predominantly adopt a document-level querying
paradigm: the entire suspicious document is encoded as a single query vector or
bag-of-words, and a relevance score is computed against each candidate
document~\cite{potthast:2014,robertson:2009}. While effective for single-source
plagiarism, this paradigm suffers from \textbf{multi-source signal dilution}.
When a suspicious text composites content from Source~A's introduction,
Source~B's methodology, and Source~C's results, the query signal becomes a noisy
mixture of multiple provenance fingerprints. Each individual source's matching
signal is diluted by noise from other sources, causing the true source document
to be ranked outside the top retrieval positions.

To address this problem, we propose \textbf{Query-Chunk Provenance Modeling}.
Rather than treating the suspicious document as a monolithic query, we decompose
it into paragraph-level semantic segments. Each segment independently retrieves
candidate sources via BM25, and a coverage-weighted voting mechanism aggregates
the results. Documents matched by multiple independent segments receive a
coverage bonus---a strong signal of genuine multi-paragraph provenance. This
approach explicitly models the provenance structure of multi-source plagiarism
while maintaining a lightweight implementation (numpy only, no GPU required).


%% ===================================================================
\section{Related Work}

Source retrieval has been a core component of PAN plagiarism detection tasks
since PAN~2014~\cite{potthast:2014}. Early approaches relied on lexical matching:
word $n$-gram overlap, fingerprinting, and BM25~\cite{robertson:2009} have
proven effective for verbatim and lightly-paraphrased plagiarism. Dense
retrieval methods using sentence encoders such as Sentence-BERT~\cite{reimers:2019}
and E5~\cite{wang:2022} project documents into a dense vector space for semantic
matching, achieving better generalization on heavily paraphrased texts.
Chunk-based retrieval~\cite{karpukhin:2020} splits documents into fixed-size
segments for independent encoding and max-pooling aggregation, but applies
splitting only on the document side while keeping the query as a single unit.
Query decomposition techniques~\cite{min:2019} split complex queries into
sub-queries for multi-hop QA, but have not been applied to plagiarism-specific
segmentation.

Our work differs from all of the above by modeling provenance structure on the
\textit{query side}: decomposing the suspicious document itself to isolate
individual source signals before retrieval. Rather than computing a single
query--document similarity, we construct an explicit provenance graph through
independent segment-level retrieval and coverage-weighted voting.


%% ===================================================================
\section{Method}

\subsection{Task Formulation}

Given a source corpus $\mathcal{C} = \{d_1, d_2, \ldots, d_N\}$
($N = 60{,}592$ academic documents from ClueWeb09) and a suspicious document
$q$, source retrieval aims to produce a ranked list
$R(q) = [d_{r_1}, \ldots, d_{r_K}]$ such that the true source document $d^*(q)$
appears as high as possible. In generative plagiarism, $q$ may be synthesized
from multiple sources, though evaluation considers one primary source per $q$.
Retrieval quality is measured by Recall@K, nDCG@K, MRR, and Reciprocal
Rank (RR).

\subsection{Query-Chunk Provenance Modeling}

Our method consists of three stages, illustrated in Figure~\ref{fig:arch}.

\begin{figure}[t]
\centering
\fbox{\begin{minipage}{0.95\linewidth}
\vspace{4em}
\begin{center}
{\small\textbf{[Stage 1]} Suspicious Document $\to$ Paragraph Split
$\to$ Query Chunks $C_1, C_2, \ldots, C_n$}\\[1em]
{\small\textbf{[Stage 2]} Each $C_i$ $\to$ Independent BM25 Retrieval
$\to$ Segment-Level Candidate Lists}\\[1em]
{\small\textbf{[Stage 3]} Coverage-Weighted Voting
$\to$ Final Ranked Source List}\\[0.5em]
{\footnotesize Dashed boxes highlight two key modules:
\underline{Paragraph Segmentation} and \underline{Coverage-Weighted Voting}.}
\vspace{4em}
\end{center}
\end{minipage}}
\caption{Overall architecture of the proposed method.}
\label{fig:arch}
\end{figure}

\subsubsection{Paragraph-Level Query Segmentation}

The suspicious document $q$ is split by double-newline paragraph boundaries.
Adjacent paragraphs are merged until reaching $\textit{max\_chars}=3{,}000$,
with chunks shorter than $\textit{min\_chars}=800$ discarded. The output is
capped at 20 chunks. This produces semantically coherent segments roughly
corresponding to academic document section structure. For documents lacking
paragraph structure, the entire text is treated as a single segment.

\subsubsection{Independent Segment Retrieval}

Each segment $c_i$ independently performs BM25 retrieval with optimized
parameters: $\textit{max\_terms}=20$ (vs.\ 100 for full-document queries, since
short segments contain fewer discriminative terms) and $\textit{max\_df}=5{,}000$
(filtering terms appearing in ${>}5{,}000$ documents, $\sim$8\% of the corpus).
Without filtering, segment-level BM25 is 85--150$\times$ slower than
full-document BM25 because segments contain mostly common academic function
words with very long posting lists (30K--50K entries). The $\textit{max\_df}$
threshold eliminates non-discriminative terms, achieving an $\sim$20$\times$
speedup (1.7--3.0\,s $\to$ 0.15\,s per segment). Heap-based top-$K$ extraction
(\texttt{heapq.nlargest}) replaces full sorting for additional efficiency.

\subsubsection{Coverage-Weighted Source Voting}

Segment-level rankings are aggregated into a final ranking via:
\begin{equation}
\text{Score}(d) = \sum_{i=1}^{n} \text{BM25}(d|c_i) \times
\left(1 + \frac{\text{count}(d)}{n}\right)
\end{equation}
where $\text{count}(d)$ is the number of segments that retrieved document $d$ in
their top-50 results, and $n$ is the total number of segments. The coverage
bonus $(1 + \text{count}(d)/n) \in [1, 2]$ rewards documents consistently
matched by multiple independent segments---a strong signal of genuine
multi-paragraph provenance rather than incidental single-phrase matching. This
mechanism can be viewed as constructing an explicit provenance graph: nodes are
query segments and candidate documents, edges represent segment--document
matches, and the final ranking reflects each document's global importance in the
graph.


%% ===================================================================
\section{Experiments}

\subsection{Experimental Setup}

\subsubsection{Dataset}
Development data comes from the PAN~2025 generative plagiarism detection dataset
converted to PAN26 format: 60,592 source documents from ClueWeb09
(avg.\ 11,765 tokens), 42,444 LLM-generated suspicious documents (training),
5,522 held-out queries with 7,949 source documents, and 43,140 query-document
relevance annotations (qrels). The PAN26 official spot-check dataset (4 queries)
was used only for format verification and TIRA submission validation. Official
main test set results are from TIRA (accessed May~25, 2026).

\subsubsection{Preprocessing}
All texts are lowercased and tokenized via regex \texttt{[a-z0-9]+}. An inverted
index is built with 1,273,004 unique terms. BM25 parameters: $k_1=1.2$,
$b=0.75$. Paragraph segmentation: $\textit{min\_chars}=800$,
$\textit{max\_chars}=3{,}000$. Segment retrieval: $\textit{max\_terms}=20$,
$\textit{max\_df}=5{,}000$, $\textit{top\_k}=50$. Standard retrieval:
$\textit{max\_terms}=100$, $\textit{top\_k}=100$. All metrics are evaluated
using \texttt{scripts/evaluate.py} (supporting multi-relevant qrels).

\subsubsection{Evaluation Metrics}
Four standard IR metrics are used: (1)~Recall@K ($K=10,100$)---proportion of
queries where the correct source appears in the top-$K$ results; (2)~nDCG@10
with position weighting; (3)~MRR---Mean Reciprocal Rank; and (4)~RR---Reciprocal
Rank, used by the official TIRA evaluation. All metrics range $[0,1]$, higher is
better.

\subsubsection{Baselines}
Three comparison categories: (1)~Standard BM25---full document as single query,
$\textit{max\_terms}=100$; (2)~E5 Chunking---256-token document chunks encoded
with E5-base-v2, coverage-aware scoring for BM25 blind-spot queries;
(3)~Ablation variants---fixed-chunk segmentation (3,000-char rolling windows)
and semantic segmentation (section header pattern matching). All methods share
identical inverted indexes and BM25 parameters for fair comparison.

\subsection{Results and Analysis}

\begin{table*}[t]
\caption{Development data method overview. *Segmented BM25 results from 800-query
sample (seed=20260520, submit\_pan26.py TIRA submission code).}
\label{tab:overview}
\begin{tabular}{lccccc}
\hline
\textbf{Method} & \textbf{R@10} & \textbf{R@100} & \textbf{nDCG@10} & \textbf{MRR} & \textbf{Data Scale} \\
\hline
Standard BM25 & 0.9196 & 0.9769 & 0.8184 & 0.7856 & 42,444q (full train) \\
+ Query Decomposition & 0.9317 & 0.9890 & 0.8274 & 0.7937 & 982q (Cat1 blind-spot) \\
+ E5 Chunking & 0.9830 & 0.9917 & 0.8716 & 0.8355 & 42,444q (full train) \\
Segmented BM25 (ours)* & 0.9700 & 0.9912 & 0.9242 & 0.9092 & 800q sample (TIRA code)\\
\hline
\end{tabular}
\end{table*}

Query Decomposition recovered 514/982 Cat1 blind-spot queries (52.3\% recovery
rate) by applying paragraph-segmented BM25 only to queries whose true source
was not retrieved by standard BM25 in the top-100. E5 Chunking attains higher
R@10 but requires a 440\,MB E5 model; our method is numpy-only.

\begin{table*}[t]
\caption{Held-out validation (7,949 documents, 5,522 queries).}
\label{tab:holdout}
\begin{tabular}{lcccc}
\hline
\textbf{Method} & \textbf{R@10} & \textbf{R@100} & \textbf{nDCG@10} & \textbf{MRR} \\
\hline
Standard BM25 & 0.9804 & 0.9874 & 0.9323 & 0.9174 \\
Segmented BM25 (ours) & 0.9947 & 0.9969 & 0.9770 & 0.9724 \\
\hline
\end{tabular}
\end{table*}

\begin{table*}[t]
\caption{Segmentation strategy comparison on the held-out set.}
\label{tab:strategies}
\begin{tabular}{lcccc}
\hline
\textbf{Strategy} & \textbf{R@10} & \textbf{MRR} & \textbf{Segmented} & \textbf{Time (s)} \\
\hline
No segmentation (BM25) & 0.9804 & 0.9174 & -- & 33 \\
Paragraph split (ours) & 0.9947 & 0.9724 & 99.9\% & 358 \\
Fixed chunk & 0.9933 & 0.9705 & 99.9\% & 300 \\
Semantic split (section headers) & 0.9746 & 0.9195 & 51.6\% & 49 \\
\hline
\end{tabular}
\end{table*}

Table~\ref{tab:strategies} reveals several findings. (1)~Paragraph segmentation
achieves the best overall performance. (2)~Fixed-chunk segmentation performs
close to paragraph split but is slightly inferior, confirming that paragraph
boundaries provide more meaningful semantic units than arbitrary character
windows. (3)~Semantic segmentation based on section header pattern matching
(Introduction, Method, Results, etc.) actually degrades performance below
standard BM25 (R@10 0.9746 vs.\ 0.9804), because only 51.6\% of queries contain
explicit section headers. This negative result is retained as a
credibility-enhancing finding: simple rule-based semantic segmentation is not
viable for this dataset.

\begin{table}[t]
\caption{Ablation study---speed optimization (200 queries).}
\label{tab:ablation}
\begin{tabular}{lcccc}
\hline
\textbf{Configuration} & \textbf{R@10} & \textbf{nDCG@10} & \textbf{MRR} & \textbf{Time/chunk} \\
\hline
(a) Standard BM25 & 0.900 & 0.805 & 0.775 & $\sim$0.02\,s \\
(b) + segmentation ($\textit{mt}{=}50$) & 0.975 & 0.942 & 0.931 & 1.7--3.0\,s \\
(c) + $\textit{max\_df}{=}5{,}000$ & 0.973 & 0.940 & 0.929 & $\sim$0.30\,s \\
(d) + $\textit{max\_terms}{=}20$ + heapq & 0.975 & 0.942 & 0.931 & $\sim$0.15\,s \\
\hline
\end{tabular}
\end{table}

Table~\ref{tab:ablation} demonstrates that the three optimization strategies
($\textit{max\_df}$ filtering, $\textit{max\_terms}$ reduction, heapq top-$K$)
together achieve an $\sim$15--20$\times$ speedup with negligible accuracy loss
($\Delta$R@10~$=$~0.002, within statistical noise).

\begin{table*}[t]
\caption{Official PAN26 main test results (TIRA, accessed May~25, 2026).}
\label{tab:tira}
\begin{tabular}{lccccc}
\hline
\textbf{System} & \textbf{Recall@10} & \textbf{Recall@100} & \textbf{nDCG@10} & \textbf{nDCG@100} & \textbf{RR} \\
\hline
BM25 baseline (official) & 0.5767 & 0.7825 & 0.5532 & 0.6068 & 0.8318 \\
fosuai / oscillating-list (ours) & 0.6300 & 0.6300 & 0.6263 & 0.6263 & 0.9350 \\
JLP / jlp (best system) & 0.8542 & 0.9629 & 0.7994 & 0.8326 & 0.9760 \\
\hline
\end{tabular}
\end{table*}

Table~\ref{tab:tira} shows our system achieves RR~$=$~0.9350 (ranked~3rd),
demonstrating strong early-rank performance consistent with the design goal
of coverage-weighted voting. Recall@10 and nDCG@10 both improve over the BM25
baseline (+9.2\% and +13.2\% respectively). However, Recall@100 equals
Recall@10 (0.6300) because the submitted run file contained only top-10
documents per query rather than the expected top-1,000. This is an engineering
limitation to be corrected in future submissions. The gap between
development-set R@10 ($\sim$0.97--0.99) and test-set R@10 (0.63) reflects
differences in corpus size, domain, and plagiarism type distribution,
underscoring that cross-dataset generalization remains the core challenge of
this task.


%% ===================================================================
\section{Conclusions}

We presented Query-Chunk Provenance Modeling, a lightweight (numpy-only) source
retrieval method for PAN26 generative plagiarism detection. The method
decomposes suspicious documents into paragraph-level segments, independently
retrieves sources per segment via BM25, and aggregates results through
coverage-weighted voting to construct an explicit provenance graph. On
development data, our method consistently outperforms standard BM25. On the
official PAN26 main test set, our system (fosuai / oscillating-list) achieved
RR~$=$~0.9350 (ranked~3rd by RR) with Recall@10 and nDCG@10 improving over the
BM25 baseline. The submitted run's top-10 output limit constrained Recall@100;
future work should output top-1,000 and explore topic-based or embedding-based
segmentation strategies for improved semantic coherence. PAN26 official
secondary dataset results are pending.


%% ===================================================================
\begin{acknowledgments}
This work is supported by the National Social Science Foundation of China
(24BYY080). We thank the PAN@CLEF~2026 organizers for providing the evaluation
dataset and the TIRA platform.
\end{acknowledgments}


%% ===================================================================
\section*{Declaration on Generative AI}

During the preparation of this work, the author(s) used Claude in order to:
English abstract translation, and English language proofreading. After using
this tool/service, the author(s) reviewed and edited the content as needed and
take(s) full responsibility for the publication's content.


%% ===================================================================
\bibliography{participant-lit,pan25-lit,PAN26_notebook}


\end{document}
